\section{Construções em Vect (Espaços Vetoriais Finitos sobre $\mathbb{C}$)}

A categoria \textbf{Vect} consiste em espaços vetoriais de dimensão finita sobre os números complexos. Os morfismos são transformações lineares. Esta categoria é completa e cocompleta.

\subsection{Fundamentação Teórica}
\begin{itemize}
    \item \textbf{Produto ($V \times W$):} Produto cartesiano dos espaços vetoriais, com a estrutura de espaço vetorial definida componente a componente.
    \item \textbf{Coproduto ($V \oplus W$):} Soma direta dos espaços vetoriais.
    \item \textbf{Equalizador:} Núcleo da transformação linear $f - g$.
    \item \textbf{Coequalizador:} Espaço quociente de $V$ pelo subespaço gerado por $\{f(v) - g(v) \mid v \in V\}$.
    \item \textbf{Pullback:} Subespaço de $V \times W$ definido por $\{(v, w) \mid f(v) = g(w)\}$.
    \item \textbf{Pushout:} Soma direta módulo a relação de equivalência gerada por $(f(v), -g(v))$ para todo $v \in U$.
\end{itemize}

\subsection{Implementação Computacional}
\begin{lstlisting}[caption={Construções Universais em Vect}]
% Produto e Projeções
% V e W são matrizes cujas linhas são vetores
prod_basis = [V; zeros(size(W, 1), size(V, 2)), W];
proj_V = [eye(size(V, 1)), zeros(size(V, 1), size(W, 1))];
proj_W = [zeros(size(W, 1), size(V, 1)), eye(size(W, 1))];

% Coproduto (Soma Direta)
coprod_basis = blkdiag(V, W);
inj_V = [eye(size(V, 1)); zeros(size(W, 1), size(V, 1))];
inj_W = [zeros(size(V, 1), size(W, 1)); eye(size(W, 1))];

% Equalizador (Núcleo de f - g)
A = f - g;
[~, ~, V_eq] = svd(A);
eq_basis = V_eq(:, size(V, 2) - rank(A) + 1:end);

% Coequalizador (Espaço Quociente)
B = [f - g];
[Q, ~] = qr(B, 'vector');
coeq_basis = Q(:, rank(B) + 1:end);

% Pullback
[U, S, V_pb] = svd([f -g]);
pb_basis = V_pb(:, size(S, 2) - rank(S) + 1:end);

% Pushout (Soma Direta módulo relação)
% Requer a definição de um espaço quociente
\end{lstlisting}